\section{Limitaciones}
A pesar de los avances en el campo de la transcripción automática de música, existen varias limitaciones que deben ser consideradas al desarrollar y evaluar sistemas de detección de acordes. En primer lugar, la calidad y diversidad del conjunto de datos utilizado para entrenar los modelos pueden afectar significativamente su rendimiento. Si el conjunto de datos no incluye una amplia variedad de estilos musicales, técnicas de ejecución y condiciones de grabación, el modelo puede tener dificultades para generalizar a situaciones del mundo real.Además, la presencia de ruido de fondo, variaciones en la afinación y técnicas de ejecución no convencionales pueden degradar el rendimiento del sistema, lo que limita su aplicabilidad en entornos no controlados~\cite{benetos2019automatic}. 

Otra limitación importante es la capacidad de los modelos para manejar la polifonía compleja de la guitarra acústica, donde múltiples notas se superponen y generan espectros armónicos complejos. Finalmente, aunque se pueden lograr altos niveles de precisión en la detección de acordes, es importante reconocer que ningún sistema de transcripción automática puede igualar completamente la capacidad de un músico profesional para interpretar y comprender la música, lo que implica que siempre habrá un margen de error y limitaciones en la interpretación de los resultados~\cite{jamshidi2024machine}.